
\documentclass{article}
\usepackage{colortbl}
\usepackage{makecell}
\usepackage{multirow}
\usepackage{supertabular}

\begin{document}

\newcounter{utterance}

\centering \large Interaction Transcript for game `dond', experiment `semi', episode 2 with mistral{-}large{-}3{-}azure.
\vspace{24pt}

{ \footnotesize  \setcounter{utterance}{1}
\setlength{\tabcolsep}{0pt}
\begin{supertabular}{c@{$\;$}|p{.15\linewidth}@{}p{.15\linewidth}p{.15\linewidth}p{.15\linewidth}p{.15\linewidth}p{.15\linewidth}}
    \# & $\;$A & \multicolumn{4}{c}{Game Master} & $\;\:$B\\
    \hline

    \theutterance \stepcounter{utterance}  
    & & \multicolumn{4}{p{0.6\linewidth}}{
        \cellcolor[rgb]{0.9,0.9,0.9}{
            \makecell[{{p{\linewidth}}}]{
                \texttt{\tiny{[P1$\langle$GM]}}
                \texttt{Jucați o jocul negocierii cooperative în care trebuie să vă puneți de acord cu un alt jucător cu privire la modul în care se împarte un set de articole între voi.} \\
\\ 
\texttt{Reguli:} \\
\texttt{(a) Tu și celălalt jucător primiți un set de articole. Fiecare dintre voi primește și o funcție de valoare secretă, reprezentând cât apreciezi fiecare tip de obiect.} \\
\texttt{(b) Schimbați mesaje cu celălalt jucător pentru a fi de acord cu privire la cine primește care articole. Fiecare poate trimite maximum 5 mesaje sau termina jocul mai devreme dând o propunere secretă în orice moment.} \\
\texttt{(c) Vi se cere să trimiteți o propunere secretă indicând articolele pe care le doriți, formatate între paranteze pătrate după cum urmează: "{[}Propunere: <număr> <nume articol>, <număr> <nume articol>, <...>{]}"} \\
\texttt{(d) Dacă propunerile tale sunt complementare, adică sunt suficiente articole pentru a îndeplini ambele propuneri, fiecare jucător primește un scor pe baza sumei valorilor articolelor pe care le{-}a primit. În caz contrar, amândoi primiți zero puncte.} \\
\texttt{(e) Scopul tău este să maximizezi suma scorului tău și a scorului celuilalt jucător.} \\
\\ 
\texttt{Să începem.} \\
\\ 
\texttt{Setul de articole disponibile este:} \\
\\ 
\texttt{1 geantă, 3 lămpi, 1 floare} \\
\\ 
\texttt{Funcția ta de valoare secretă este:} \\
\\ 
\texttt{geantă: 3, lampă: 0, floare: 7} \\
\\ 
\texttt{IMPORTANT: Mesajele tale, cu excepția unei propuneri secrete, sunt transmise direct celuilalt jucător, deci nu include niciun răspuns la reguli sau text care anunță mesajul tău. Pentru a face o propunere secretă, utilizează formatul indicat. Nu utiliza paranteze pătrate când comunici cu celălalt jucător, altfel va fi interpretat ca propunerea ta secretă.} \\
            }
        }
    }
    & \\ \\

    \theutterance \stepcounter{utterance}  
    & \multicolumn{4}{p{0.6\linewidth}}{
        \cellcolor[rgb]{0.9,0.9,0.9}{
            \makecell[{{p{\linewidth}}}]{
                \texttt{\tiny{[P1$\rangle$GM]}}
                \texttt{Bună! Ce părere ai să începem prin a discuta ce articole ne interesează cel mai mult? Eu sunt interesat în special de floare. Tu ce preferi?} \\
            }
        }
    }
    & & \\ \\

    \theutterance \stepcounter{utterance}  
    & & & \multicolumn{2}{p{0.3\linewidth}}{
        \cellcolor[rgb]{0.9,0.9,0.9}{
            \makecell[{{p{\linewidth}}}]{
                \texttt{\tiny{[GM$|$GM]}}
                \texttt{continue} \\
            }
        }
    }
    & & \\ \\

    \theutterance \stepcounter{utterance}  
    & & & \multicolumn{2}{p{0.3\linewidth}}{
        \cellcolor[rgb]{0.9,0.9,0.9}{
            \makecell[{{p{\linewidth}}}]{
                \texttt{\tiny{[GM$|$GM]}}
                \texttt{Bună! Ce părere ai să începem prin a discuta ce articole ne interesează cel mai mult? Eu sunt interesat în special de floare. Tu ce preferi?} \\
            }
        }
    }
    & & \\ \\

    \theutterance \stepcounter{utterance}  
    & & \multicolumn{4}{p{0.6\linewidth}}{
        \cellcolor[rgb]{0.9,0.9,0.9}{
            \makecell[{{p{\linewidth}}}]{
                \texttt{\tiny{[GM$\rangle$P2]}}
                \texttt{Jucați o jocul negocierii cooperative în care trebuie să vă puneți de acord cu un alt jucător cu privire la modul în care se împarte un set de articole între voi.} \\
\\ 
\texttt{Reguli:} \\
\texttt{(a) Tu și celălalt jucător primiți un set de articole. Fiecare dintre voi primește și o funcție de valoare secretă, reprezentând cât apreciezi fiecare tip de obiect.} \\
\texttt{(b) Schimbați mesaje cu celălalt jucător pentru a fi de acord cu privire la cine primește care articole. Fiecare poate trimite maximum 5 mesaje sau termina jocul mai devreme dând o propunere secretă în orice moment.} \\
\texttt{(c) Vi se cere să trimiteți o propunere secretă indicând articolele pe care le doriți, formatate între paranteze pătrate după cum urmează: "{[}Propunere: <număr> <nume articol>, <număr> <nume articol>, <...>{]}"} \\
\texttt{(d) Dacă propunerile tale sunt complementare, adică sunt suficiente articole pentru a îndeplini ambele propuneri, fiecare jucător primește un scor pe baza sumei valorilor articolelor pe care le{-}a primit. În caz contrar, amândoi primiți zero puncte.} \\
\texttt{(e) Scopul tău este să maximizezi suma scorului tău și a scorului celuilalt jucător.} \\
\\ 
\texttt{Să începem.} \\
\\ 
\texttt{Setul de articole disponibile este:} \\
\\ 
\texttt{1 geantă, 3 lămpi, 1 floare} \\
\\ 
\texttt{Funcția ta de valoare secretă este:} \\
\\ 
\texttt{geantă: 1, lampă: 2, floare: 3} \\
\\ 
\texttt{IMPORTANT: Mesajele tale, cu excepția unei propuneri secrete, sunt transmise direct celuilalt jucător, deci nu include niciun răspuns la reguli sau text care anunță mesajul tău. Pentru a face o propunere secretă, utilizează formatul indicat. Nu utiliza paranteze pătrate când comunici cu celălalt jucător, altfel va fi interpretat ca propunerea ta secretă.} \\
\\ 
\texttt{Bună! Ce părere ai să începem prin a discuta ce articole ne interesează cel mai mult? Eu sunt interesat în special de floare. Tu ce preferi?} \\
            }
        }
    }
    & \\ \\

    \theutterance \stepcounter{utterance}  
    & & & \multicolumn{4}{p{0.6\linewidth}}{
        \cellcolor[rgb]{0.9,0.9,0.9}{
            \makecell[{{p{\linewidth}}}]{
                \texttt{\tiny{[GM$\langle$P2]}}
                \texttt{Ce zici să împărțim articolele astfel: tu iei floarea și o lampă, iar eu iau geanta și celelalte două lămpi? Sau ai o altă idee?} \\
            }
        }
    }
     \\ \\

    \theutterance \stepcounter{utterance}  
    & & & \multicolumn{2}{p{0.3\linewidth}}{
        \cellcolor[rgb]{0.9,0.9,0.9}{
            \makecell[{{p{\linewidth}}}]{
                \texttt{\tiny{[GM$|$GM]}}
                \texttt{continue} \\
            }
        }
    }
    & & \\ \\

    \theutterance \stepcounter{utterance}  
    & & & \multicolumn{2}{p{0.3\linewidth}}{
        \cellcolor[rgb]{0.9,0.9,0.9}{
            \makecell[{{p{\linewidth}}}]{
                \texttt{\tiny{[GM$|$GM]}}
                \texttt{Ce zici să împărțim articolele astfel: tu iei floarea și o lampă, iar eu iau geanta și celelalte două lămpi? Sau ai o altă idee?} \\
            }
        }
    }
    & & \\ \\

    \theutterance \stepcounter{utterance}  
    & & \multicolumn{4}{p{0.6\linewidth}}{
        \cellcolor[rgb]{0.9,0.9,0.9}{
            \makecell[{{p{\linewidth}}}]{
                \texttt{\tiny{[P1$\langle$GM]}}
                \texttt{Ce zici să împărțim articolele astfel: tu iei floarea și o lampă, iar eu iau geanta și celelalte două lămpi? Sau ai o altă idee?} \\
            }
        }
    }
    & \\ \\

    \theutterance \stepcounter{utterance}  
    & \multicolumn{4}{p{0.6\linewidth}}{
        \cellcolor[rgb]{0.9,0.9,0.9}{
            \makecell[{{p{\linewidth}}}]{
                \texttt{\tiny{[P1$\rangle$GM]}}
                \texttt{Mulțumesc pentru propunere! Eu aș prefera să iau doar floarea, deoarece pentru mine are cea mai mare valoare. Ce zici să îmi lași mie floarea și să iei tu geanta și toate cele 3 lămpi? Astfel, amândoi am putea obține un scor bun.} \\
            }
        }
    }
    & & \\ \\

    \theutterance \stepcounter{utterance}  
    & & & \multicolumn{2}{p{0.3\linewidth}}{
        \cellcolor[rgb]{0.9,0.9,0.9}{
            \makecell[{{p{\linewidth}}}]{
                \texttt{\tiny{[GM$|$GM]}}
                \texttt{continue} \\
            }
        }
    }
    & & \\ \\

    \theutterance \stepcounter{utterance}  
    & & & \multicolumn{2}{p{0.3\linewidth}}{
        \cellcolor[rgb]{0.9,0.9,0.9}{
            \makecell[{{p{\linewidth}}}]{
                \texttt{\tiny{[GM$|$GM]}}
                \texttt{Mulțumesc pentru propunere! Eu aș prefera să iau doar floarea, deoarece pentru mine are cea mai mare valoare. Ce zici să îmi lași mie floarea și să iei tu geanta și toate cele 3 lămpi? Astfel, amândoi am putea obține un scor bun.} \\
            }
        }
    }
    & & \\ \\

    \theutterance \stepcounter{utterance}  
    & & \multicolumn{4}{p{0.6\linewidth}}{
        \cellcolor[rgb]{0.9,0.9,0.9}{
            \makecell[{{p{\linewidth}}}]{
                \texttt{\tiny{[GM$\rangle$P2]}}
                \texttt{Mulțumesc pentru propunere! Eu aș prefera să iau doar floarea, deoarece pentru mine are cea mai mare valoare. Ce zici să îmi lași mie floarea și să iei tu geanta și toate cele 3 lămpi? Astfel, amândoi am putea obține un scor bun.} \\
            }
        }
    }
    & \\ \\

    \theutterance \stepcounter{utterance}  
    & & & \multicolumn{4}{p{0.6\linewidth}}{
        \cellcolor[rgb]{0.9,0.9,0.9}{
            \makecell[{{p{\linewidth}}}]{
                \texttt{\tiny{[GM$\langle$P2]}}
                \texttt{Îmi place ideea de a maximiza scorurile amândurora. Dacă eu iau geanta și toate cele 3 lămpi, iar tu iei floarea, scorul meu ar fi 1 (geantă) + 3 x 2 (lămpi) = 7, iar al tău ar fi 3 (floare). Totalul ar fi 10, ceea ce este foarte bun.} \\
\\ 
\texttt{Dar, pentru a fi siguri că suntem amândoi mulțumiți, ce zici să facem următoarea propunere secretă:} \\
\\ 
\texttt{{[}Propunere: 1 floare{]}} \\
            }
        }
    }
     \\ \\

    \theutterance \stepcounter{utterance}  
    & & & \multicolumn{2}{p{0.3\linewidth}}{
        \cellcolor[rgb]{0.9,0.9,0.9}{
            \makecell[{{p{\linewidth}}}]{
                \texttt{\tiny{[GM$|$GM]}}
                \texttt{proposal} \\
            }
        }
    }
    & & \\ \\

    \theutterance \stepcounter{utterance}  
    & & & \multicolumn{2}{p{0.3\linewidth}}{
        \cellcolor[rgb]{0.9,0.9,0.9}{
            \makecell[{{p{\linewidth}}}]{
                \texttt{\tiny{[GM$|$GM]}}
                [0, 0, 1]
            }
        }
    }
    & & \\ \\

    \theutterance \stepcounter{utterance}  
    & & \multicolumn{4}{p{0.6\linewidth}}{
        \cellcolor[rgb]{0.9,0.9,0.9}{
            \makecell[{{p{\linewidth}}}]{
                \texttt{\tiny{[P1$\langle$GM]}}
                \texttt{Celălalt jucător a trimis o propunere secretă. Acum este rândul tău să trimiți o propunere proprie.} \\
            }
        }
    }
    & \\ \\

    \theutterance \stepcounter{utterance}  
    & \multicolumn{4}{p{0.6\linewidth}}{
        \cellcolor[rgb]{0.9,0.9,0.9}{
            \makecell[{{p{\linewidth}}}]{
                \texttt{\tiny{[P1$\rangle$GM]}}
                \texttt{{[}Propunere: 1 floare{]}} \\
            }
        }
    }
    & & \\ \\

    \theutterance \stepcounter{utterance}  
    & & & \multicolumn{2}{p{0.3\linewidth}}{
        \cellcolor[rgb]{0.9,0.9,0.9}{
            \makecell[{{p{\linewidth}}}]{
                \texttt{\tiny{[GM$|$GM]}}
                \texttt{proposal} \\
            }
        }
    }
    & & \\ \\

    \theutterance \stepcounter{utterance}  
    & & & \multicolumn{2}{p{0.3\linewidth}}{
        \cellcolor[rgb]{0.9,0.9,0.9}{
            \makecell[{{p{\linewidth}}}]{
                \texttt{\tiny{[GM$|$GM]}}
                [0, 0, 1]
            }
        }
    }
    & & \\ \\

    \theutterance \stepcounter{utterance}  
    & & & \multicolumn{2}{p{0.3\linewidth}}{
        \cellcolor[rgb]{0.9,0.9,0.9}{
            \makecell[{{p{\linewidth}}}]{
                \texttt{\tiny{[GM$|$GM]}}
                [[0, 0, 1], [0, 0, 1]]
            }
        }
    }
    & & \\ \\

\end{supertabular}
}

\end{document}
